\documentclass[10pt,a4paper]{amsart}
\usepackage[margin=19mm]{geometry}
\usepackage{amsmath,amssymb,mathtools}
\usepackage[scr=rsfs,cal=euler]{mathalfa}
\usepackage{xfrac}
\usepackage[unicode,hidelinks,bookmarks=false]{hyperref}
\newcommand{\ie}{i.e.\ }
\newcommand{\cf}{cf.\ }
\newcommand{\resp}{respectively\ }
\newcommand{\Subsection}{\subsection}
\providecommand{\nobreakdash}{\nobreak\hbox{-}}
\setlength{\emergencystretch}{2em}
\multlinegap0pt
\usepackage{mathtools}
\usepackage{mathrsfs}
\usepackage{booktabs,array}
\usepackage{lmodern}
\setlength{\parindent}{0pt}
\setlength{\parskip}{0.45em}
\allowdisplaybreaks
\newtheorem{theorem}{Theorem}[section]
\newtheorem{lemma}[theorem]{Lemma}
\newtheorem{proposition}[theorem]{Proposition}
\newtheorem{corollary}[theorem]{Corollary}
\newtheorem{definition}[theorem]{Definition}
\newtheorem{remark}[theorem]{Remark}
\newtheorem{example}[theorem]{Example}
\newtheorem{conjecture}[theorem]{Conjecture}
\newcommand{\Energy}{\mathcal E}
\begin{document}
\title[Collision Positivity for Three-Variable Symmetric Monomial I]{Collision Positivity for Three-Variable Symmetric Monomial Inequalities}
\author{Jian Sun}
\date{}
\maketitle
\noindent\emph{Source and licence.} Collision Positivity for Three-Variable Symmetric Monomial Inequalities. Version arXiv:2609.22465v1 (CC BY 4.0). This material is distributed under the Creative Commons Attribution 4.0 International licence, http://creativecommons.org/licenses/by/4.0/, as recorded in the arXiv OAI licence metadata for this exact version and shown on the abstract page https://arxiv.org/abs/2609.22465v1. Full rights notice: licenses/arxiv-2609.22465-CC-BY-4.0.txt.\par\smallskip
\noindent\textit{Editorial scope.} Counted unit: the original \texttt{theorem} environment \texttt{thm:minimal-state-reduction} at source lines 1232--1259 together with its complete proof at lines 1261--1292; the delivered document keeps source lines 1013--1293 so that every referenced label and every citation resolves inside it. Native editable source is the paper's own concatenated TeX; the AMS wrapper is editorial formatting only. No publisher class, upstream script or unrelated artwork is redistributed.
\section{Lattice Structure and Reduction to Minimal Collision-Positive States}
\label{sec:monge}

We now vary the middle partition and the lower endpoint while keeping the upper endpoint
\begin{equation*}
\lambda=(p,q,0)
\end{equation*}
fixed.  For integers \(r,s,\beta,\kappa\), set
\begin{equation*}
\gamma_{r,s}
=
(p-r-s,\ q+r,\ s),
\end{equation*}
and
\begin{equation*}
\mu_\kappa
=
(p-\beta,\ q+\beta-\kappa,\ \kappa).
\end{equation*}
For every state for which these triples are admissible, define
\begin{equation*}
P_{r,s,\kappa}
=
J_\lambda+J_{\mu_\kappa}-2J_{\gamma_{r,s}},
\end{equation*}
and define the collision quotient \(Q_{r,s,\kappa}\) by
\begin{equation*}
P_{r,s,\kappa}(t,1,1)
=
2(t-1)^2Q_{r,s,\kappa}(t).
\end{equation*}
The purpose of this section is to identify the monotone lattice directions of
\(P_{r,s,\kappa}\) and \(Q_{r,s,\kappa}\), and then to reduce the proof to states that are minimal with respect to those directions.

\subsection{Legal states and finite-difference identities}

\begin{definition}[Legal state]
\label{def:legal-state}
A triple \((r,s,\kappa)\) is called \emph{legal}, for fixed \((p,q,\beta)\), if
\begin{equation*}
\lambda\succcurlyeq\gamma_{r,s}\succcurlyeq\mu_\kappa
\end{equation*}
and both \(\gamma_{r,s}\) and \(\mu_\kappa\) are nonnegative weakly decreasing triples.
Equivalently,
\begin{align*}
s&\ge0,
& r+s&\ge0,
& q+r-s&\ge0,
& d-2r-s&\ge0,\\
\beta-r-s&\ge0,
& \kappa-s&\ge0,
& d-2\beta+\kappa&\ge0,
& q+\beta-2\kappa&\ge0.
\end{align*}
A neighboring state is called legal if the shifted triple satisfies the same conditions.
\end{definition}

The first four inequalities state that \(\gamma_{r,s}\) is a partition dominated by \(\lambda\); the next two are exactly the two partial-sum inequalities for
\(\gamma_{r,s}\succcurlyeq\mu_\kappa\); and the last two state that \(\mu_\kappa\) is weakly decreasing.  Notice in particular that legality is a property of the lattice state itself, not an additional hypothesis introduced in the finite-difference formulas below.

For any lattice quantity \(F_{r,s,\kappa}\), define the forward difference operators by
\begin{align*}
\Delta_rF_{r,s,\kappa}
&=
F_{r+1,s,\kappa}-F_{r,s,\kappa},\\
\Delta_sF_{r,s,\kappa}
&=
F_{r,s+1,\kappa}-F_{r,s,\kappa},\\
\Delta_\kappa F_{r,s,\kappa}
&=
F_{r,s,\kappa+1}-F_{r,s,\kappa}.
\end{align*}
Thus
\begin{equation*}
\Delta_r\Delta_sF_{r,s,\kappa}
=
F_{r+1,s+1,\kappa}
-F_{r+1,s,\kappa}
-F_{r,s+1,\kappa}
+F_{r,s,\kappa}.
\end{equation*}
The positive monotonicity direction for the lower endpoint is \(-\Delta_\kappa\), because increasing \(\kappa\) moves \(\mu_\kappa\) downward in dominance order.

\begin{lemma}[Finite-difference monotonicity identities]
\label{lem:finite-difference-monotonicity}
Let \((r,s,\kappa)\) be legal.
Whenever the indicated neighboring state is also legal, one has
\begin{align*}
\Delta_rP_{r,s,\kappa}
&=
2\bigl(J_{\gamma_{r,s}}-J_{\gamma_{r+1,s}}\bigr)
\ge0,\\
\Delta_sP_{r,s,\kappa}
&=
2\bigl(J_{\gamma_{r,s}}-J_{\gamma_{r,s+1}}\bigr)
\ge0,
\end{align*}
and
\begin{equation*}
-\Delta_\kappa P_{r,s,\kappa}
=
J_{\mu_\kappa}-J_{\mu_{\kappa+1}}
\ge0.
\end{equation*}
At the collision quotient level,
\begin{align*}
\Delta_rQ_{r,s,\kappa}(t)
&=
2t^{q+r}A_{d-2r-s-1}(t),\\
\Delta_sQ_{r,s,\kappa}(t)
&=
2t^sA_{p-r-2s-1}(t),\\
-\Delta_\kappa Q_{r,s,\kappa}(t)
&=
t^\kappa A_{q+\beta-2\kappa-1}(t).
\end{align*}
Moreover, whenever the corresponding two-dimensional configuration is legal,
\begin{equation*}
\Delta_r\Delta_sQ_{r,s,\kappa}(t)
=
-2t^{p-r-s-2},
\end{equation*}
whereas the lower-endpoint direction is separable from the two middle-partition directions:
\begin{equation*}
\Delta_\kappa\Delta_rQ_{r,s,\kappa}
=
\Delta_\kappa\Delta_sQ_{r,s,\kappa}
=
0.
\end{equation*}
\end{lemma}

\begin{proof}
The three full-polynomial inequalities follow from Muirhead's inequality.  Indeed,
\(\gamma_{r,s}\) dominates \(\gamma_{r+1,s}\) by one unit transfer from the first coordinate to the second,
\(\gamma_{r,s}\) dominates \(\gamma_{r,s+1}\) by one unit transfer from the first coordinate to the third, and
\(\mu_\kappa\) dominates \(\mu_{\kappa+1}\) by one unit transfer from the second coordinate to the third.

Because the orbit sums are labeled, for any exponent triple \((u,v,w)\),
\begin{equation*}
J_{(u,v,w)}(t,1,1)
=
2(t^u+t^v+t^w).
\end{equation*}
For integers \(u>v\),
\begin{align*}
t^u+t^v-t^{u-1}-t^{v+1}
&=(t-1)(t^{u-1}-t^v)\\
&=(t-1)^2t^vA_{u-v-1}(t).
\end{align*}
Applying this identity to the three unit transfers above and dividing by
\(2(t-1)^2\) gives the three formulas for \(Q\).

For the mixed difference, use the first collision finite-difference formula:
\begin{align*}
\Delta_s\Delta_rQ_{r,s,\kappa}(t)
&=
2t^{q+r}
\left(
A_{d-2r-s-2}(t)-A_{d-2r-s-1}(t)
\right)\\
&=
-2t^{q+r}t^{d-2r-s-2}\\
&=
-2t^{p-r-s-2}.
\end{align*}
Finally, the formulas for \(\Delta_rQ\) and \(\Delta_sQ\) do not involve \(\kappa\), which gives
\begin{equation*}
\Delta_\kappa\Delta_rQ
=
\Delta_\kappa\Delta_sQ
=0.
\end{equation*}
\end{proof}

For each fixed \(t>0\), Lemma~\ref{lem:finite-difference-monotonicity} shows that
\(Q_{r,s,\kappa}(t)\) is nondecreasing in \(r\) and \(s\), and nonincreasing in \(\kappa\), along legal neighboring states.  In addition,
\begin{equation*}
\Delta_r\Delta_sQ_{r,s,\kappa}(t)<0.
\end{equation*}
We shall refer to this sign condition as the \emph{discrete Monge property} of the \((r,s)\)-array.  The vanishing mixed differences with \(\kappa\) mean that the lower-endpoint increment is independent of the two middle-partition increments.  These formulas, rather than any separate geometric slogan, are the lattice structure used in the reduction below.

\subsection{Reduction to minimal collision-positive states}

We write
\begin{equation*}
Q_{r,s,\kappa}\ge0
\end{equation*}
to mean
\begin{equation*}
Q_{r,s,\kappa}(t)\ge0
\qquad\text{for every }t>0.
\end{equation*}
Similarly,
\begin{equation*}
Q_{r,s,\kappa}\not\ge0
\end{equation*}
means that \(Q_{r,s,\kappa}(t)<0\) for at least one \(t>0\).

\begin{definition}[Predecessors and minimal collision-positive states]
\label{def:minimal-collision-positive}
For a legal state \((r,s,\kappa)\), its legal predecessors are the states
\begin{equation*}
(r-1,s,\kappa),
\qquad
(r,s-1,\kappa),
\qquad
(r,s,\kappa+1),
\end{equation*}
whenever they are legal.
A legal state \((r,s,\kappa)\) is called \emph{minimal collision-positive} if
\begin{equation*}
Q_{r,s,\kappa}\ge0
\end{equation*}
and every legal predecessor fails collision positivity.
\end{definition}

The terminology is consistent with the monotonicity directions in Lemma~\ref{lem:finite-difference-monotonicity}: moving from a predecessor to \((r,s,\kappa)\) adds a globally nonnegative full-polynomial increment.


\begin{theorem}[Minimal-State Reduction]
\label{thm:minimal-state-reduction}
Let \((r,s,\kappa)\) be a collision-positive legal state.  Then there exists a minimal collision-positive legal state
\begin{equation*}
(r_*,s_*,\kappa_*)
\end{equation*}
with
\begin{equation*}
r_*\le r,
\qquad
s_*\le s,
\qquad
\kappa_*\ge\kappa,
\end{equation*}
such that
\begin{equation*}
P_{r,s,\kappa}(x,y,z)
\ge
P_{r_*,s_*,\kappa_*}(x,y,z)
\qquad(x,y,z>0).
\end{equation*}
Consequently, if
\begin{equation*}
P_{r_*,s_*,\kappa_*}(x,y,z)\ge0
\qquad(x,y,z>0)
\end{equation*}
for every minimal collision-positive legal state, then the same conclusion holds for every collision-positive legal state.
\end{theorem}

\begin{proof}
Start from a collision-positive legal state \((r,s,\kappa)\).  If it is not minimal collision-positive, at least one of its legal predecessors is collision-positive.  Move to such a predecessor and repeat.

The legal lattice is finite: with the total degree and the upper endpoint fixed, only finitely many integer partitions can occur as the middle and lower exponent vectors.  In addition, every predecessor step strictly decreases the integer
\begin{equation*}
r+s-\kappa,
\end{equation*}
so no cycle is possible.  Hence the descent terminates at a minimal collision-positive legal state \((r_*,s_*,\kappa_*)\).

At every step Lemma~\ref{lem:finite-difference-monotonicity} gives a globally nonnegative difference.  For example,
\begin{align*}
P_{r,s,\kappa}-P_{r-1,s,\kappa}
&=
\Delta_rP_{r-1,s,\kappa}
\ge0,\\
P_{r,s,\kappa}-P_{r,s-1,\kappa}
&=
\Delta_sP_{r,s-1,\kappa}
\ge0,\\
P_{r,s,\kappa}-P_{r,s,\kappa+1}
&=
-\Delta_\kappa P_{r,s,\kappa}
\ge0.
\end{align*}
Telescoping along the finite predecessor chain therefore gives
\begin{equation*}
P_{r,s,\kappa}
\ge
P_{r_*,s_*,\kappa_*}
\end{equation*}
pointwise on \((0,\infty)^3\).  The final assertion follows immediately.
\end{proof}

\end{document}
